\documentclass[12pt]{article}
\usepackage[utf8]{inputenc}
\usepackage[T1]{fontenc}
\usepackage[margin=1in]{geometry}
\usepackage{setspace}
\usepackage{xcolor}
\usepackage[colorlinks=true, urlcolor=blue!60!black, linkcolor=blue!60!black]{hyperref}
\setlength{\parindent}{0em}
\setlength{\parskip}{0.6em}
\renewcommand{\baselinestretch}{1.0}
\begin{document}
\begin{center}
{\LARGE \textit{Marsh v.\ Alabama}, 326 U.S.\ 501 (1946)\par}
\vspace{0.5em}
{\large Case Brief --- Two-Page Edition\par}
\end{center}

\section{Memory Jogger}
A company town performs a public function and is bound by the First and Fourteenth Amendments, so Alabama could not convict a Jehovah's Witness of trespass for leafletting on its sidewalk (\href{https://supreme.justia.com/cases/federal/us/326/501/}{Justia}).

\section{Detailed Case Facts}
Chickasaw, Alabama, was wholly owned by the Gulf Shipbuilding Corporation but had every attribute of a town: homes, streets, sewers, a business block, a company-paid deputy sheriff, and a U.S.\ post office, all open to highway travelers. Grace Marsh, a Jehovah's Witness, distributed religious literature on the business-block sidewalk despite a company notice barring solicitation, and refused to leave.

\section{Procedural History}
Marsh was convicted under Alabama's criminal trespass statute; the Alabama Court of Appeals affirmed because Gulf had not dedicated the sidewalk to the public, and the state supreme court denied review. The U.S.\ Supreme Court reversed.

\section{Judicial Votes}
Reversed 5--3. Black wrote for Douglas, Murphy, and Rutledge, with Frankfurter joining in part and concurring; Reed (with Stone and Burton) dissented; Jackson did not participate.

\section{Holding}
A State may not criminally punish a person for distributing religious literature on the freely accessible sidewalk of a company-owned town. A corporation that governs a community with all the attributes of a municipality performs a ``public function,'' and its property rights do not override First Amendment liberties, which occupy a ``preferred position.''

\section{Analysis of Opinions}
\textbf{Majority:} an ordinary town could not ban the leafletting; state-law title is not dispositive; facilities run for the public perform ``essentially a public function''; and residents of company towns are free citizens who need uncensored information. \textbf{Frankfurter:} ``a company-owned town is a town,'' but he rejected the Commerce Clause analogy. \textbf{Reed:} the ``novel Constitutional doctrine'' overrides property rights, and Marsh could have leafletted on the public highway thirty feet away.

\section{Examples: Future Applications}
\textbf{Same-Side:} a corporate ``tech town'' with housing, utilities, a post office, and deputized security. \textbf{Opposite-Side:} a dominant social-media platform, which merely hosts speech (\textit{Halleck}). \textbf{Fence-Sitter:} a developer-owned town whose police, fire, and schools are public.

\section{Critique}
\textit{Marsh} blends public-function and state-enforcement theories and frames the result as balancing, while modern doctrine treats state action as a threshold test confined to functions ``traditionally exclusively'' public.

\section{Key Quotations}
\begin{quote}
``Ownership does not always mean absolute dominion. The more an owner, for his advantage, opens up his property for use by the public in general, the more do his rights become circumscribed by the statutory and constitutional rights of those who use it.''
\end{quote}
\noindent --- Justice Black, for the Court, 326 U.S.\ at 506.

\end{document}
